\documentclass[11pt]{book}
\usepackage[paperwidth=145mm,paperheight=200mm,top=18mm,bottom=20mm,inner=20mm,outer=16mm]{geometry}
\usepackage{brumfiel}
% figures15.tex -- Chapter 15 (Analytic Geometry) figures redrawn in TikZ.
% Coordinate-geometry chapter: every figure is exactly constructible from the
% mathematics in the text (axes, plotted points, lines, circles), so vertices
% are computed, never eyeballed. Proportions were measured off the printed
% figures at 400 dpi (PDF pp. 267-279) during the figure review; the readings
% are recorded in chapters/ch15-UNCERTAIN.md.
%
% Helper macros used only inside this file (check_labels greps FIG[A-Z]+ for
% the figure macros, so these helpers are not picked up as figures).
% `positioning' is needed for the "above left=<a> and <b>" node syntax; it is
% loaded here rather than in brumfiel.sty because this file is also compiled
% standalone by tools/check_labels.py.
\usetikzlibrary{positioning}

% Standard axes WITH the letters x and y: {xmin}{xmax}{ymin}{ymax}.
% The book letters the axes in Figs 15-6 to 15-19 and leaves them unlettered
% in Fig 15-4, so the two variants are kept apart.
\newcommand{\XVaxes}[4]{%
  \draw[fig,-latex] (#1,0) -- (#2,0);
  \node[vlab,right] at (#2,0) {$x$};
  \draw[fig,-latex] (0,#3) -- (0,#4);
  \node[vlab,right] at (0,#4) {$y$};}
% Same axes, unlettered (Fig 15-4).
\newcommand{\XVaxesnl}[4]{%
  \draw[fig,-latex] (#1,0) -- (#2,0);
  \draw[fig,-latex] (0,#3) -- (0,#4);}
% Origin label
\newcommand{\XVzero}{\node[vlab,below left=-1pt and 0pt] at (0,0) {$0$};}
% Unit ticks on the x-axis, from #1 to #2 step 1 (skipping 0)
\newcommand{\XVxticks}[2]{%
  \foreach \t in {#1,...,#2}{\ifnum\t=0\else
    \draw[tick] (\t,-0.11) -- (\t,0.11);\fi}}
\newcommand{\XVyticks}[2]{%
  \foreach \t in {#1,...,#2}{\ifnum\t=0\else
    \draw[tick] (-0.11,\t) -- (0.11,\t);\fi}}

% ---------------------------------------------------------------- Fig 15-1
% Coordinatized line: O, the unit point P_1, a positive point P_x and a
% negative point P'_{x'}.  Collinear by construction.  Spacings follow the
% printed figure (O -> P_1 : P_1 -> P_x : P'_{x'} -> O = 1 : 1.55 : 1.22).
\newcommand{\FIGXVONE}{\begin{bkfigure}[\linewidth]
\begin{tikzpicture}[scale=1.0]
  \draw[fig] (-0.3,0) -- (10.3,0);
  \foreach \x/\below/\above in {0.6/{P'_{x'}}/{x'}, 3.5/{O}/{},
                                5.9/{P_1}/{1}, 9.6/{P_x}/{x}}
    {\dt{\x,0}
     \node[vlab,below=2pt] at (\x,0) {$\below$};}
  \foreach \x/\above in {0.6/{x'}, 5.9/{1}, 9.6/{x}}
    {\node[vlab,above=2pt] at (\x,0) {$\above$};}
\end{tikzpicture}
\figcap{15--1}
\end{bkfigure}}

% ---------------------------------------------------------------- Fig 15-2
% Two PERPENDICULAR lines through O, coordinatized, drawn obliquely to make
% the point that the choice of axes is arbitrary.
% abscissa positive direction 210 deg; ordinate positive direction 300 deg.
% Measured off the plate: abscissa runs 2.6 units back / 5.0 forward, ordinate
% 1.5 back / 3.4 forward.  Every numeral sits on the side the book puts it on:
% the positive abscissa numerals on the 120 deg side, the positive ordinate
% numerals on the 30 deg side, and each -1 thrown to the far side of its own
% axis so the two do not crowd the 0.
\newcommand{\FIGXVTWO}{\begin{bkfigure}[\linewidth]
\begin{tikzpicture}[scale=1.040]
  % unit vectors
  \def\ax{-0.86603} \def\ay{-0.5}      % abscissa positive (210 deg)
  \def\ox{0.5}      \def\oy{-0.86603}  % ordinate positive (300 deg)
  % the two axes
  \draw[fig] (2.6*-\ax,2.6*-\ay) -- (5.0*\ax,5.0*\ay);
  \draw[fig] (1.5*-\ox,1.5*-\oy) -- (3.4*\ox,3.4*\oy);
  % ticks + numerals on the abscissa axis (tick direction = ordinate dir)
  \foreach \t/\n in {1/1, 2/2, 3/3}
    {\draw[tick] ({\t*\ax-0.20*\ox},{\t*\ay-0.20*\oy})
              -- ({\t*\ax+0.20*\ox},{\t*\ay+0.20*\oy});
     \node[slab] at ({\t*\ax-0.46*\ox},{\t*\ay-0.46*\oy}) {$\n$};}
  \draw[tick] ({-\ax-0.20*\ox},{-\ay-0.20*\oy})
           -- ({-\ax+0.20*\ox},{-\ay+0.20*\oy});
  \node[slab] at ({-\ax-0.50*\ox},{-\ay-0.50*\oy}) {$-1$};
  % ticks + numerals on the ordinate axis (tick direction = abscissa dir)
  \foreach \t/\n in {1/1, 2/2, 3/3}
    {\draw[tick] ({\t*\ox-0.20*\ax},{\t*\oy-0.20*\ay})
              -- ({\t*\ox+0.20*\ax},{\t*\oy+0.20*\ay});
     \node[slab] at ({\t*\ox-0.50*\ax},{\t*\oy-0.50*\ay}) {$\n$};}
  \draw[tick] ({-\ox-0.20*\ax},{-\oy-0.20*\ay})
           -- ({-\ox+0.20*\ax},{-\oy+0.20*\ay});
  \node[slab] at ({-\ox+0.66*\ax},{-\oy+0.66*\ay}) {$-1$};
  % 0 in the open sector between the two negative rays, as in the book
  \node[slab] at (0.16,0.60) {$0$};
  \node[slab,left]  at (5.15*\ax,5.15*\ay) {Abscissas};
  \node[slab,right] at (3.55*\ox,3.55*\oy) {Ordinates};
\end{tikzpicture}
\figcap{15--2}
\end{bkfigure}}

% ---------------------------------------------------------------- Fig 15-3
% Theorem 15-1: P,Q in the SAME quadrant => PQ meets neither axis;
% P,Q in DIFFERENT quadrants => PQ meets at least one axis.
\newcommand{\FIGXVTHREE}{\begin{bkfigure}[\linewidth]
\begin{tikzpicture}[scale=1.014]
  \draw[fig] (-4.3,0) -- (3.9,0);
  \draw[fig] (0,2.9) -- (0,-1.5);
  % same quadrant (I): segment meets neither axis
  \coordinate (P1) at (1.6,2.15);  \coordinate (Q1) at (2.65,1.35);
  \draw[key] (P1) -- (Q1);
  \node[vlab,above=1pt] at (P1) {$P$};
  \node[vlab,right=1pt] at (Q1) {$Q$};
  % quadrant II to quadrant III: crosses the axis of abscissas
  \coordinate (Q2) at (-2.45,1.5); \coordinate (P2) at (-3.75,-0.62);
  \draw[key] (Q2) -- (P2);
  \node[vlab,above=1pt] at (Q2) {$Q$};
  \node[vlab,left=1pt]  at (P2) {$P$};
  % quadrant II to quadrant IV: crosses both axes
  \coordinate (P3) at (-1.25,1.05); \coordinate (Q3) at (2.35,-0.72);
  \draw[key] (P3) -- (Q3);
  \node[vlab,above=1pt] at (P3) {$P$};
  \node[vlab,right=1pt] at (Q3) {$Q$};
  \foreach \p in {P1,Q1,Q2,P2,P3,Q3} {\dt{\p}}
\end{tikzpicture}
\figcap{15--3}
\end{bkfigure}}

% ---------------------------------------------------------------- Fig 15-4
% The four quadrants, numbered.  The book leaves this pair of axes unlettered
% (no x, no y) and runs the left branch longer than the right.
\newcommand{\FIGXVFOUR}{\begin{bkfigure}[\linewidth]
\begin{tikzpicture}[scale=0.936]
  \XVaxesnl{-3.6}{2.9}{-1.2}{2.7}
  \XVxticks{1}{2}
  \XVyticks{1}{2}
  \node[slab,above] at (1,0.12) {$1$};
  \node[slab,above] at (2,0.12) {$2$};
  \node[slab,left]  at (-0.12,1) {$1$};
  \node[slab,left]  at (-0.12,2) {$2$};
  \XVzero
  \node[vlab] at (1.40,1.50)  {\textup{I}};
  \node[vlab] at (-1.70,1.50) {\textup{II}};
  \node[vlab] at (-1.75,-0.87){\textup{III}};
  \node[vlab] at (1.36,-0.87) {\textup{IV}};
\end{tikzpicture}
\figcap{15--4}
\end{bkfigure}}

% ---------------------------------------------------------------- Fig 15-5
% Rotated (still perpendicular) axes; the reader is asked which quadrant is
% which.  Ordinate positive direction 150 deg, abscissa positive 240 deg --
% the printed pair measures about 92 deg apart, so they are squared up.
% Branch lengths follow the plate: the ordinate axis runs 1.8 units forward
% and 3.9 back, the abscissa 2.05 forward and 1.95 back.
\newcommand{\FIGXVFIVE}{\begin{bkfigure}[\linewidth]
\begin{tikzpicture}[scale=1.040]
  \def\ox{-0.86603} \def\oy{0.5}       % ordinate positive (150 deg)
  \def\ax{-0.5}     \def\ay{-0.86603}  % abscissa positive (240 deg)
  \draw[fig,-latex] (3.9*-\ox,3.9*-\oy) -- (1.8*\ox,1.8*\oy);
  \draw[fig,-latex] (1.95*-\ax,1.95*-\ay) -- (2.05*\ax,2.05*\ay);
  \dt{\ox,\oy}
  \dt{\ax,\ay}
  \node[slab,above] at (\ox,\oy) {$1$};
  \node[slab,right] at (\ax,\ay) {$1$};
  \node[slab,above] at (2.0*\ox,2.0*\oy) {Ordinates};
  \node[slab,below] at (2.25*\ax,2.25*\ay) {Abscissas};
\end{tikzpicture}
\figcap{15--5}
\end{bkfigure}}

% ---------------------------------------------------------------- Fig 15-6
% Definition 15-4: l_1 through P parallel to the axis of ordinates meets the
% first axis in A; l_2 through P parallel to the axis of abscissas meets the
% second axis in B.  A sits on the positive abscissa axis and B on the
% negative ordinate axis, so P falls in quadrant IV, as in the book.
\newcommand{\FIGXVSIX}{\begin{bkfigure}[\linewidth]
\begin{tikzpicture}[scale=1.040]
  \XVaxes{-1.6}{2.8}{-3.2}{1.65}
  \coordinate (A) at (1.55,0);
  \coordinate (B) at (0,-1.70);
  \coordinate (P) at (1.55,-1.70);
  \draw[aux] (1.55,1.55) -- (1.55,-3.5);
  \draw[aux] (-1.65,-1.70) -- (3.2,-1.70);
  \foreach \p in {A,B,P} {\dt{\p}}
  \node[vlab,above left=-1pt and 0pt] at (0,0) {$O$};
  \node[vlab,above right=1pt and 1pt] at (A) {$A$};
  \node[vlab,above left=1pt and 1pt]  at (B) {$B$};
  \node[vlab,above right=1pt and 1pt] at (P) {$P$};
  \node[vlab,right] at (3.2,-1.70) {$l_2$};
  \node[vlab,below right=-1pt and 0pt] at (1.55,-3.5) {$l_1$};
\end{tikzpicture}
\figcap{15--6}
\end{bkfigure}}

% ---------------------------------------------------------------- Fig 15-7
% The line x = 3: parallel to the y-axis, 3 units to its right.  The book
% ticks only 1 and 2 -- the line itself stands in for the third mark -- and
% sets the label to the RIGHT of the line.
\newcommand{\FIGXVSEVEN}{\begin{bkfigure}[\linewidth]
\begin{tikzpicture}[scale=0.910]
  \XVaxes{-0.9}{4.3}{-1.7}{3.55}
  \XVxticks{1}{2}
  \draw[key] (3,-1.65) -- (3,3.48);
  \XVzero
  \node[vlab,right] at (3.08,1.2) {$x=3$};
\end{tikzpicture}
\figcap{15--7}
\end{bkfigure}}

% ------------------------------------------------------------ Fig 15-8, 15-9
% 15-8: the points of 2x - 3y + 6 = 0 found in Example 1, namely
%       (0,2), (3,4), (-3,0) -- collinear by construction.  Both axes are
%       ticked in the book.
% 15-9: the line y = -2 of Example 2, parallel to the x-axis, 2 units below.
%       The axis of ordinates stops on that line, as it does in the book.
\newcommand{\FIGXVEIGHTNINE}{\begin{bkfigure}[\linewidth]
\begin{minipage}{0.48\linewidth}\centering
\begin{tikzpicture}[scale=0.44]
  \XVaxes{-4.6}{6.4}{-2.6}{5.4}
  \XVxticks{-4}{5}
  \XVyticks{-2}{4}
  % the line 2x - 3y + 6 = 0, i.e. y = (2/3)x + 2 -- endpoints computed
  \draw[aux] (-4.2,-0.8) -- (5.4,5.6);
  \foreach \p in {(-3,0),(0,2),(3,4)} {\node[bkdot] at \p {};}
\end{tikzpicture}
\figcap{15--8}
\end{minipage}\hfill
\begin{minipage}{0.48\linewidth}\centering
\begin{tikzpicture}[scale=0.44]
  \XVaxes{-3.4}{5.3}{-2.05}{3.1}
  \draw[aux] (-3.6,-2) -- (5.4,-2);
  \draw[tick] (-0.22,-1) -- (0.22,-1);
  \XVzero
  \node[vlab,above=1pt] at (3.15,-2) {$(x,-2)$};
\end{tikzpicture}
\figcap{15--9}
\end{minipage}
\end{bkfigure}}

% ---------------------------------------------------------- Fig 15-10, 15-11
% 15-10: proof of Theorem 15-2, Case 3.  (0,-C/B), (x1,y1), (x2,y2) are
%        collinear; the two vertical legs are perpendicular to the base.
%        The base is taken at height 0 and the x-axis sits below it, at the
%        same fraction of x2 as in the book.
% 15-11: proof of Theorem 15-3, Case 3.  (x1,y1), (x2,y2), (x,y) collinear;
%        right angles at (x2,y1) and (x,y1).  In the book (x1,y1) lies just
%        to the LEFT of the axis of ordinates and the whole construction sits
%        above the axis of abscissas.
\newcommand{\FIGXVTENELEVEN}{\begin{bkfigure}[\linewidth]
\begin{minipage}{0.49\linewidth}\centering
\begin{tikzpicture}[scale=0.55]
  % base line y = -C/B taken at height 0; the x-axis sits below it.
  \def\bY{0}
  \coordinate (O)  at (0,\bY);          % (0, -C/B)
  \coordinate (P1) at (2.81,1.47525);   % (x1, y1)   slope 0.525
  \coordinate (P2) at (4.2,2.205);      % (x2, y2)   -- collinear with O,P1
  \coordinate (F1) at (2.81,\bY);
  \coordinate (F2) at (4.2,\bY);
  \draw[fig,-latex] (-0.52,-1.65) -- (4.54,-1.65);
  \node[vlab,right] at (4.54,-1.65) {$x$};
  \draw[fig,-latex] (0,-2.20) -- (0,2.49);
  \node[vlab,right] at (0,2.49) {$y$};
  \draw[fig] (O) -- (F2);
  \draw[fig] (P1) -- (F1);
  \draw[fig] (P2) -- (F2);
  \draw[aux] (O) -- (P2);
  \foreach \p in {O,P1,P2} {\dt{\p}}
  \node[slab,left]  at (O)  {$(0,-C/B)$};
  \node[slab,anchor=south east,xshift=-1pt,yshift=3pt] at (P1) {$(x_1,y_1)$};
  \node[slab,anchor=south east,xshift=9pt,yshift=3pt]  at (P2) {$(x_2,y_2)$};
  % one node holds both foot labels, as the book sets them: side by side on a
  % single line below the base.  Two separate nodes overlap at this trim.
  \node[slab,below] at (3.3,\bY) {$(x_1,-C/B)\quad(x_2,-C/B)$};
\end{tikzpicture}
\figcap{15--10}
\end{minipage}\hfill
\begin{minipage}{0.49\linewidth}\centering
\begin{tikzpicture}[scale=0.60]
  \coordinate (P1) at (-0.61,1.36);     % (x1, y1)
  \coordinate (P2) at (1.78,1.77825);   % (x2, y2)   slope 0.175
  \coordinate (P3) at (4.67,2.284);     % (x, y)   -- collinear with P1,P2
  \coordinate (F2) at (1.78,1.36);
  \coordinate (F3) at (4.67,1.36);
  \draw[fig,-latex] (-1.65,0) -- (4.93,0);
  \node[vlab,right] at (4.93,0) {$x$};
  \draw[fig,-latex] (0,-1.02) -- (0,3.44);
  \node[vlab,right] at (0,3.44) {$y$};
  \draw[fig] (P1) -- (F3);
  \draw[fig] (P2) -- (F2);
  \draw[fig] (P3) -- (F3);
  \draw[fig] (P1) -- (P3);
  \foreach \p in {P1,P2,P3} {\dt{\p}}
  \node[slab,left] at (P1) {$(x_1,y_1)$};
  \node[slab] at ($(P2)+(-0.34,0.48)$) {$(x_2,y_2)$};
  \node[slab] at ($(P3)+(0.10,0.50)$)  {$(x,y)$};
  \node[slab,below] at (F2) {$(x_2,y_1)$};
  \node[slab,below] at (F3) {$(x,y_1)$};
\end{tikzpicture}
\figcap{15--11}
\end{minipage}
\end{bkfigure}}

% --------------------------------------------------------------- Fig 15-12
% The line through (-2,3) and (1,5) of the worked Example.  Both axes ticked;
% the book sets (-2,3) above-left of its point and (1,5) below-right of its
% own, which also keeps both clear of the line.
\newcommand{\FIGXVTWELVE}{\begin{bkfigure}[\linewidth]
\begin{tikzpicture}[scale=0.676]
  \XVaxes{-5.0}{5.6}{-1.2}{6.9}
  \XVxticks{-4}{5}
  \XVyticks{1}{5}
  % slope 2/3 through (-2,3) and (1,5) -- endpoints computed, not eyeballed
  \draw[key] (-4.1,1.6) -- (3.45,6.63333);
  \dt{-2,3}
  \dt{1,5}
  \XVzero
  \node[vlab,above left=0pt and 0pt]  at (-2,3) {$(-2,3)$};
  \node[vlab,below right=0pt and 0pt] at (1,5)  {$(1,5)$};
\end{tikzpicture}
\figcap{15--12}
\end{bkfigure}}

% ---------------------------------------------------------- Fig 15-13, 15-14
% 15-13: distances along lines parallel to an axis.  AB = 1 - (-3) = 4;
%        CD = |x2 - x1|.  All four dashed pairs are axis-parallel.  C and D
%        straddle the axis of ordinates, at (-1,-1) and (2,-1), as in the
%        book, and each of the four lettered points carries its letter on one
%        side of the dot and its coordinate pair on the other.
% 15-14: the right triangle P1 P2 Q of the distance formula, right angle at Q.
\newcommand{\FIGXVTHIRTEENFOURTEEN}{\begin{bkfigure}[\linewidth]
\begin{minipage}{0.58\linewidth}\centering
\begin{tikzpicture}[scale=0.52]
  \draw[fig,-latex] (-2.8,0) -- (7.0,0);
  \node[vlab,right] at (7.0,0) {$x$};
  \draw[fig,-latex] (0,-3.6) -- (0,4.6);
  \node[vlab,right] at (0,4.6) {$y$};
  \node[slab,above left=-1pt and 0pt] at (0,0) {$0$};
  % horizontal pair (2,3)-(6,3)
  \draw[aux] (2,3) -- (6,3);
  \dt{2,3} \dt{6,3}
  \node[slab,above] at (2,3.1) {$(2,3)$};
  \node[slab,above] at (6,3.1) {$(6,3)$};
  % horizontal pair (-2,2)-(3,2)
  \draw[aux] (-2,2) -- (3,2);
  \dt{-2,2} \dt{3,2}
  \node[slab,left]  at (-2.1,2) {$(-2,2)$};
  \node[slab,right] at (3.1,2)  {$(3,2)$};
  % vertical pair A(1,1)-B(1,-3)
  \draw[aux] (1,1) -- (1,-3);
  \dt{1,1} \dt{1,-3}
  \node[slab,left]  at (0.88,1)  {$A$};
  \node[slab,right] at (1.12,1)  {$(1,1)$};
  \node[slab,left]  at (0.88,-3) {$B$};
  \node[slab,right] at (1.12,-3) {$(1,-3)$};
  % horizontal pair C(x1,y1)-D(x2,y1), straddling the axis of ordinates.
  % Dropped to y = -1.5: at this trim a label set above a point at y = -1
  % would run back over the axis of abscissas.
  \draw[aux] (-1,-1.5) -- (2,-1.5);
  \dt{-1,-1.5} \dt{2,-1.5}
  \node[slab,above] at (-1,-1.4) {$C$};
  \node[slab,above] at (2,-1.4)  {$D$};
  \node[slab,left]  at (-1.12,-1.5) {$(x_1,y_1)$};
  \node[slab,right] at (2.12,-1.5)  {$(x_2,y_1)$};
\end{tikzpicture}
\figcap{15--13}
\end{minipage}\hfill
\begin{minipage}{0.38\linewidth}\centering
\begin{tikzpicture}[scale=0.62]
  \draw[fig,-latex] (-0.95,0) -- (4.57,0);
  \node[vlab,right] at (4.57,0) {$x$};
  \draw[fig,-latex] (0,-0.26) -- (0,3.55);
  \node[vlab,right] at (0,3.55) {$y$};
  \node[slab,below right=-1pt and 1pt] at (0,0) {$0$};
  \coordinate (P1) at (0.62,1.15);
  \coordinate (P2) at (3.72,3.05);
  \coordinate (Q)  at (3.72,1.15);
  \draw[key] (P1) -- (P2);
  \draw[aux] (P1) -- (Q);
  \draw[aux] (Q)  -- (P2);
  \foreach \p in {P1,P2,Q} {\dt{\p}}
  \node[vlab,above,xshift=-2pt] at (P1) {$P_1$};
  \node[vlab,anchor=south east,xshift=-1pt] at (P2) {$P_2$};
  \node[vlab,right,yshift=5pt] at (Q) {$Q$};
  % nudged right so it clears the axis of ordinates, as the book's does
  \node[slab,below,xshift=8pt] at (P1) {$(x_1,y_1)$};
  % the book prints (x_2,y_1) at P_2 as well as at Q -- see ch15-UNCERTAIN.md
  \node[slab,anchor=south west,xshift=1pt] at (P2) {$(x_2,y_1)$};
  \node[slab,below=3pt] at (Q) {$(x_2,y_1)$};
\end{tikzpicture}
\figcap{15--14}
\end{minipage}
\end{bkfigure}}

% ---------------------------------------------------------- Fig 15-15, 15-16
% 15-15: M is the midpoint of P1P2; parallels to the axes through M bisect
%        the sides of the right triangle.
% 15-16: the worked example -- M(0,1) is the midpoint of (-1,3) and (1,-1).
\newcommand{\FIGXVFIFTEENSIXTEEN}{\begin{bkfigure}[\linewidth]
\begin{minipage}{0.52\linewidth}\centering
\begin{tikzpicture}[scale=0.756]
  \draw[fig,-latex] (-0.40,0) -- (3.60,0);
  \node[vlab,right] at (3.60,0) {$x$};
  \draw[fig,-latex] (0,-0.35) -- (0,3.90);
  \node[vlab,right] at (0,3.90) {$y$};
  \node[slab,below right=-1pt and 1pt] at (0,0) {$0$};
  \coordinate (P1) at (0.46,1.05);
  \coordinate (P2) at (3.36,2.52);
  \coordinate (M)  at ($(P1)!0.5!(P2)$);
  \coordinate (C)  at (3.36,1.05);
  \coordinate (Mh) at (3.36,1.785);    % foot of the parallel to the x-axis
  \coordinate (Mv) at (1.91,1.05);     % foot of the parallel to the y-axis
  \draw[key] (P1) -- (P2);
  \draw[aux] (P1) -- (C) -- (P2);
  \draw[aux] (M) -- (Mh);
  \draw[aux] (M) -- (Mv);
  \foreach \p in {P1,P2,M} {\dt{\p}}
  \node[slab,anchor=north west,xshift=-9pt,yshift=-3pt] at (P1) {$P_1\,(x_1,y_1)$};
  \node[slab,right] at (P2) {$P_2\,(x_2,y_2)$};
  \node[vlab,above left=-3pt and -1pt] at (M) {$M$};
  % Solidus form: a built-up fraction draws a rule inside the label box, which
  % the collision audit reads as geometry.  The solidus runs about two and a
  % half times as wide as the book's stacked fraction, so the label is lifted
  % clear of P_2 rather than tucked beside it.  See chapters/ch15-UNCERTAIN.md.
  \node[slab,anchor=south west] at (0.05,2.95)
    {$\left((x_1+x_2)/2,\ (y_1+y_2)/2\right)$};
\end{tikzpicture}
\figcap{15--15}
\end{minipage}\hfill
\begin{minipage}{0.44\linewidth}\centering
\begin{tikzpicture}[scale=0.651]
  \XVaxes{-2.0}{3.2}{-1.95}{3.5}
  \node[slab,below left=-1pt and 0pt] at (0,0) {$0$};
  \dt{-1,3}
  \dt{0,1}
  \dt{1,-1}
  \node[slab,left]  at (-1.15,3) {$P_1\,(-1,3)$};
  \node[slab,right] at (0.18,1)  {$M\,(0,1)$};
  \node[slab,right] at (1.18,-1) {$P_2\,(1,-1)$};
\end{tikzpicture}
\figcap{15--16}
\end{minipage}
\end{bkfigure}}

% --------------------------------------------------------------- Fig 15-17
% Exercise *7: the segment joining the midpoints of two sides is parallel to
% the third side and half as long.  Vertices (0,0), (a,0), (b,c); the book
% draws b barely larger than a, so the right side is all but vertical, and
% the midline is a plain dashed segment with no arrowhead.
\newcommand{\FIGXVSEVENTEEN}{\begin{bkfigure}[\linewidth]
\begin{tikzpicture}[scale=0.936]
  \draw[fig,-latex] (-0.70,0) -- (3.70,0);
  \node[vlab,right] at (3.70,0) {$x$};
  \draw[fig,-latex] (0,-0.60) -- (0,2.75);
  \node[vlab,right] at (0,2.75) {$y$};
  \coordinate (A) at (0,0);
  \coordinate (B) at (3.10,0);
  \coordinate (C) at (3.15,2.48);
  \draw[fig] (A) -- (B) -- (C) -- cycle;
  \coordinate (M1) at ($(A)!0.5!(C)$);
  \coordinate (M2) at ($(B)!0.5!(C)$);
  \draw[aux] (M1) -- (M2);
  \foreach \p in {M1,M2} {\dt{\p}}
  \node[slab,below right=-1pt and 0pt] at (A) {$(0,0)$};
  \node[slab,below] at (B) {$(a,0)$};
  \node[slab,right] at (C) {$(b,c)$};
\end{tikzpicture}
\figcap{15--17}
\end{bkfigure}}

% ---------------------------------------------------------- Fig 15-18, 15-19
% 15-18: the circle of Eq. (15-6): centre (h,k), radius r, a point (x,y) on it.
%        The book's centre sits just above the axis of abscissas and about
%        half a radius right of the axis of ordinates, so the circle crosses
%        both axes; the axis of ordinates stops well short of the bottom of
%        the circle.
% 15-19: the circle x^2+y^2+6x-8y+19 = 0, centre (-3,4), radius sqrt(6).
\newcommand{\FIGXVEIGHTEENNINETEEN}{\begin{bkfigure}[\linewidth]
\begin{minipage}{0.50\linewidth}\centering
\begin{tikzpicture}[scale=0.66]
  \draw[fig,-latex] (-1.82,0) -- (4.90,0);
  \node[vlab,right] at (4.90,0) {$x$};
  \draw[fig,-latex] (0,-0.77) -- (0,2.65);
  \node[vlab,right] at (0,2.65) {$y$};
  % the book sets 0 below left of the origin; the origin here lies inside the
  % circle, and the audit approximates each quarter-arc by its chord, so a
  % label dropped below left reads as touching that chord.  Set above left.
  \node[slab,above left=-1pt and 0pt] at (0,0) {$0$};
  \coordinate (K) at (1.04,0.38);
  \def\rr{1.85}
  \draw[key] (K) circle (\rr);
  % (x,y) placed on the circle at 148.5 deg from the centre -- exact by
  % construction, so the drawn radius really has length r.
  \coordinate (X) at ($(K)+(148.5:\rr)$);
  \draw[fig] (K) -- (X);
  \dt{K}
  \node[slab,right] at (K) {$(h,k)$};
  \node[slab,left=1pt]  at (X) {$(x,y)$};
  % r set off the radius along the perpendicular, so it does not sit on the line
  \node[slab] at ($(K)!0.5!(X)+(58.5:0.34)$) {$r$};
\end{tikzpicture}
\figcap{15--18}
\end{minipage}\hfill
\begin{minipage}{0.46\linewidth}\centering
\begin{tikzpicture}[scale=0.48]
  \draw[fig,-latex] (-6.0,0) -- (3.2,0);
  \node[vlab,right] at (3.2,0) {$x$};
  \draw[fig,-latex] (0,-1.0) -- (0,7.0);
  \node[vlab,right] at (0,7.0) {$y$};
  \node[slab,below right=-1pt and 1pt] at (0,0) {$0$};
  \XVxticks{-5}{-1}
  \XVyticks{1}{5}
  \draw[key] (-3,4) circle (2.4494897);
  \dt{-3,4}
  \node[slab,below=3pt] at (-3,4) {$(-3,4)$};
\end{tikzpicture}
\figcap{15--19}
\end{minipage}
\end{bkfigure}}


% ---- local macros (hoist into book preamble) ----
% Starred (optional) theorems, as the book marks Theorems 15-2 and 15-3.
% Shares the theorem counter so the numbering runs 15-1, 15-2, 15-3.
\theoremstyle{bkplain}
\newtheorem{XVstheorem}[theorem]{*Theorem}
% Worked examples, set as the book sets them: italic lead-in, no number when
% a section carries only one.
\newenvironment{XVexample}[1][]{\par\medskip\noindent\emph{Example#1.}\ }{\par\medskip}
% NUMBERED remarks. brumfiel.sty's `remark' hard-codes an unnumbered
% "Remark."; book p.261 prints "Remark 1." and "Remark 2." in sequence, so the
% numbered form needs its own environment. The unnumbered `remark' is still
% used for the two lone remarks on pp.257 and 262.
\newenvironment{XVremark}[1]{\par\medskip\noindent\emph{Remark #1.}\ }{\par\medskip}
% Symbol-marked footnotes. The book marks this chapter's footnotes with * and
% dagger, never with numbers; \footnotetext would print the footnote counter,
% so the counter is blanked and the symbol is typed into the body text.
% (Same device as ch02's \IIstarnote and ch09's \IXstarnote.)
\newcommand{\XVnote}[1]{\begingroup\renewcommand{\thefootnote}{}%
  \footnotetext{#1}\endgroup}
% (No star-convention footnote here: ch15 has starred exercises and two starred
% theorems, but none of book pp.251-265 carries the "*Starred exercises ... are
% optional" note; ch02:311 owns the book-wide first use.)
% ---- end local macros ----

\begin{document}
\setcounter{chapter}{15}

\chapter*{\raggedright\Huge ANALYTIC GEOMETRY}
\addcontentsline{toc}{chapter}{15.\ Analytic Geometry}
\markboth{ANALYTIC GEOMETRY}{}
\vspace{-1em}\noindent\rule{\linewidth}{0.6pt}\vspace{1em}
\figurekey

%========================================================= 15-1
\section*{15--1\quad SOME HISTORY}
\addcontentsline{toc}{section}{15--1\quad Some history}

We have commented before on the fact that numbers, other than integers, did
not occur explicitly in Greek geometry. One reason for this was that numbers
were hard to work with at that time. Even after the introduction, during the
Middle Ages, of Arabic numerals and the place value of the decimal system,
geometry continued in the Greek spirit. As a result, it often happened that
problems that we would handle naturally with algebra were solved by geometric
constructions.

Of course, as time passed, there was a gradual tendency toward the use of
algebra in geometric problems, but there was nothing systematic in its use.
And until there was sufficient algebraic machinery available and people became
skilled in using numbers, there was little advantage in changing a geometric
problem into an algebraic problem.

By the beginning of the 17th century, considerable progress had been made in
algebra and in algebraic notation;* in particular, quadratic, cubic, and
quartic equations could be solved. Mathematicians were interested in studying
tangents to curves, which is more easily done with a little analytic geometry.
So it was that in the 17th century the times were right for something new.
Indeed, many men had actually been using some analytic geometry without making
it a definite system for attacking problems in geometry. Today, looking back
on those times, it is hard to imagine why some ideas were so slow in coming.
But we must remember that discovery is always much harder than understanding
what has been discovered.%
\XVnote{*The beginning student hardly realizes how much easier the shorthand
notation of algebra makes his problems. The interested student is urged to
read in the history of mathematics and to try working a problem in earlier
notations.}

It is universally agreed that the French philosopher Ren\'e Descartes
(1596--1650) was the discoverer (inventor?) of analytic geometry, and the date
is generally set as 1637, when his \emph{Geometry} was published. Descartes is
also noted for his work in physics and in philosophy. In the latter, he
aspired to write a philosophy that would be perfectly logical and free from
objectionable assumptions. His efforts in philosophy include the book
\emph{Discourse on Method}, which expounds what is referred to today as
Cartesian philosophy, and which has been a source of dispute ever since! But
his \emph{Geometry} is undisputed. Analytic geometry is today a part of a
liberal education as well as a necessity for pure and applied science.

Some brief remarks about Descartes' life may be of interest. Born in 1596, he
started school in a Jesuit institution at the age of eight. Never a strong,
healthy boy, he was allowed by his school to stay in bed mornings to do his
studies, and later in life he did much of his writing in bed. He left school
in 1612, studied law, and even spent some time in the Dutch army as a
gentleman volunteer, leaving in 1619. It appears that by 1620 most of his
great ideas had been established in outline. In other words, he had thought
them over and had a fairly clear picture of how they should go. Much of his
life thereafter was spent in Holland. A devout man, he nevertheless had
troubles with the church over his philosophical ideas. By the time his
\emph{Geometry} was published, he was a well-known figure. In 1649 he visited
the court of Queen Christina of Sweden, after much urging, to teach her his
system of philosophy. Assigned the outlandish hour of 5 a.m. three days a week
for the instructions, he fell ill in February 1650 and died a few days later.

\emph{Geometry} was his only book on mathematics and, in a sense, contained
nothing totally new. What Descartes did was to show how a systematic use of
coordinates helps to attack problems. He classified various kinds of curves
and suggested further work. The book was suggestive and tremendously fruitful.
We owe to him our convention of using the first letters of the alphabet
($a$, $b$, $c$) for knowns and the last letters ($x$, $y$, $z$) for unknowns.

It is not easy to give a concise definition of analytic geometry. In a sense,
this entire chapter is such a definition, for a complete definition should
include some of its uses. But for the moment, it may suffice to say that, in
analytic geometry, geometric figures are studied by means of equations. Then,
if we wish to know something concerning the geometric figure, it \emph{may} be
easy to learn this by examining an equation. The equations that are used
involve numbers called coordinates, which serve to locate points. But why the
word ``analytic''? Analytic geometry does enable one to analyze geometric
figures, but this is also done in regular Euclidean geometry. Perhaps a better
description would be to call this new way of looking at geometry, plane
\emph{coordinate} geometry, or plane \emph{Cartesian} geometry. Remember that
it is the same plane we are studying all the time, no matter how we look at it.

\exercisegroup{15--1}
\begin{multicols}{2}
\begin{exlist}
\item Read the history of mathematics during the 17th century.
\item Read a biography of Descartes in an encyclopedia or in \emph{Men of
  Mathematics}, by E. T. Bell.
\end{exlist}
\end{multicols}

%========================================================= 15-2
\section*{15--2\quad EXAMPLES}
\addcontentsline{toc}{section}{15--2\quad Examples}

Although analytic geometry is distinguished by use of coordinates, the idea
should already be familiar from some common situations. Some cities are laid
out in squares, with the east-west and north-south roads called avenues and
streets, respectively. The numbers of street and avenue fix any intersection.
In the same way, you unconsciously use coordinates in describing a nearby town
as, perhaps, 7 miles north and 5 miles west of ``here.''

\exercisegroup{15--2}
\begin{multicols}{2}
\begin{exlist}
\item Try to think of other examples of coordinates, that is situations in
  which position is determined by a pair of numbers.
\item Discuss the need to have streets straight in order to ``coordinatize'' a
  city.
\end{exlist}
\end{multicols}

By far the most important example of coordinates is our method (Chapter 5) of
putting coordinates on a line. We recall how this was done:

Given a line $l$ and a unit of length, choose a point $O$ on the line $l$, as
in Fig.~15--1. Now choose* a point $P_1$ on $l$ at a distance 1 from $O$. This
choice determines which half of the line we will call the \emph{positive
axis}, or \emph{positive direction}. For every positive real number $x$, there
is exactly one point $P_x$ of $l$ on the same side of $O$ as the point $P_1$
such that the length $\sg{OP_x} = x$. For every negative real number $x'$,
there is exactly one point $P_{x'}$ of $l$ on the opposite side of $O$ from
$P_1$ such that the length $\sg{OP_{x'}} = -x' = |x'|$.%
\XVnote{*There are two ways of choosing this point, namely on either side of
$O$. Which side is chosen is not important.}

Each of the numbers $x$ and $x'$ is called the \emph{coordinate} of the
corresponding point of the line.

\FIGXVONE

\exercisegroup{15--3}
\begin{multicols}{2}
\begin{exlist}
\item How are the points $P_{x_1}$, $P_{x_2}$, and $O$ related if $x_1 < 0$
  and $x_2 > 0$? if $0 < x_2 < x_1$? if $x_1 < x_2 < 0$?
\item If we agree that points with larger coordinates ``lie to the right'' of
  points with lesser coordinates, how are we thinking of the drawn line $l$
  and the chosen point $P_1$?
\item Two points on the line have coordinates $-3$ and 7. What is the distance
  between them? between 5 and 12? between $x_1$ and $x_2$?
\end{exlist}
\end{multicols}

\begin{definition}
We will refer to the above construction of coordinates as a coordinatization
of the line $l$, with the point $O$ as \emph{origin of coordinates}.
\end{definition}

%========================================================= 15-3
\section*{15--3\quad THE EXISTENCE OF COORDINATES}
\addcontentsline{toc}{section}{15--3\quad The existence of coordinates}

The construction of plane coordinates is best carried out in a number of
simple steps. We assume that a unit of length has been selected in the plane.

\begin{definition}
Choose a point $O$ in the plane. Call this point the \emph{origin of
coordinates}. Choose any two perpendicular lines through $O$. Call these lines
\emph{coordinate axes}. Coordinatize each of the axes, with the point $O$ as
origin of coordinates. Select one of these lines and call it the \emph{axis of
abscissas}.* The other axis is called the \emph{axis of ordinates}.\dag\ See
Fig.~15--2.%
\XVnote{*This is not a fortunately chosen word. It would be better to say
``first coordinates.'' In most drawings in textbooks, the axis of abscissas is
drawn horizontally as the reader views it. There is nothing sacred about this,
we just agree to consider one coordinate before we do the other (something
must be first) and call this first coordinate by the strange name
\emph{abscissa}.}%
\XVnote{\dag\ This word is not much better, but the literature of mathematics
is filled with this word too, so make the best of it and learn it. Of course
``second coordinate'' is much better.}
\end{definition}

\FIGXVTWO

\exercisegroup{15--4}
\begin{multicols}{2}
\begin{exlist}
\item Draw a figure, different from Fig.~15--2, illustrating Definition
  15--2.
\item Compare your figure with your neighbor's. Did you both choose the
  ``same'' axis for abscissas? Turn your figure sideways. Do you have a
  different opinion about which is the first coordinate?
\item With the same drawing as in Exercise 1, reverse the coordinates on each
  of the axes. Is this a possible way of satisfying Definition 15--2? Suppose
  that we reversed the coordinates on only one of the axes.
\end{exlist}
\end{multicols}

Perhaps you will have observed a certain arbitrariness in the definition, in
that we \emph{choose} a point $O$ and \emph{choose} lines for axes. This is
actually the way things operate in practice. In applying analytic geometry to
problems in physics or engineering, it is usually necessary to \emph{choose}
coordinates. \emph{There are no fixed coordinate systems inherent in nature.}
Coordinate systems are convenient inventions that facilitate mathematical
computations.

Having set up our axes, we can make certain statements concerning them. We
give some of these in a theorem.

\begin{theorem}
The coordinate axes separate the plane into four regions, called quadrants,
such that if $P$ and $Q$ are points not on either axis and in the same
quadrant, then the segment $PQ$ does not meet either axis, whereas if $P$ and
$Q$ are in different quadrants, then $PQ$ meets at least one of the axes. See
Fig.~15--3.
\end{theorem}

\FIGXVTHREE

\noindent\emph{Proof.}
Each axis separates the plane. (Why?) Consider any point $P$ not on any axis.
The point $P$ is on one side of the axis of abscissas and one side of the axis
of ordinates. Let $Q$ be any point on the same side of the axis of abscissas
as $P$ and on the same side of the axis of ordinates as $P$. Then $PQ$ meets
neither axis. (Why?) The set of all such points $Q$ is a quadrant. If $Q_1$
and $Q_2$ are two points of that quadrant, then $Q_1Q_2$ does not intersect
either axis. (Why?)

There are four ways of choosing the point $P$, according to which sides of the
axes it is on. This gives rise to the four quadrants.

\exercisegroup{15--5}
\begin{multicols}{2}
\begin{exlist}
\item Describe the four ways of choosing $P$ to get the four quadrants.
\item Does the above proof require the axes to be perpendicular?
\end{exlist}
\end{multicols}

\begin{definition}
The half plane that contains the positive axis of abscissas will be called the
\emph{right half plane}, and similarly for \emph{left half plane}.* The half
plane that contains the positive axis of ordinates will be called the
\emph{upper half plane}, and similarly for \emph{lower half plane}.%
\XVnote{*This agrees with the way axes are conventionally drawn in textbooks.}
\end{definition}

\FIGXVFOUR

The quadrant that is in the right half plane and in the upper half plane will
be called the \emph{first quadrant;} the \emph{second quadrant} lies in the
upper and left half planes; the \emph{third quadrant}, in the lower and left
half planes; and the \emph{fourth quadrant}, in the lower and right half
planes. See Fig.~15--4.

\exercisegroup{15--6}
\begin{multicols}{2}
\begin{exlist}
\item In Fig.~15--5, which is quadrant I? quadrant II? quadrant III? quadrant
  IV?
\item[\stex 2.] Explain why Definition 15--3 does not require any previous
  knowledge of what the words ``left'' and ``right'' mean.
\item[\stex 3.] Has there been any use made of the fact that the axes have
  been selected to be perpendicular? What changes would have to be made so far
  if the axes were not at right angles?
\end{exlist}
\end{multicols}

\FIGXVFIVE

\begin{definition}
Let $P$ be any point of the plane, and $l_1$ and $l_2$ the lines through $P$
parallel, respectively, to the axis of ordinates and axis of abscissas. The
line $l_1$ meets the first axis in a point $A$, with coordinate $x$. The line
$l_2$ meets the second axis in a point $B$, with coordinate $y$. See
Fig.~15--6. The two numbers $x$ and $y$ are called the \emph{coordinates} of
$P$. We write them in parentheses, $(x, y)$, and refer to $(x, y)$ as a
\emph{point}. The number $x$ is the \emph{abscissa} of $P$, and the number $y$
is the \emph{ordinate} of $P$.
\end{definition}

Because the letters $x$ and $y$ are so frequently used for abscissa and
ordinate, the axes are often called the \emph{x-axis} and the \emph{y-axis}.
The plane then, when coordinatized in this way, is often called the
\emph{xy-plane}.\dag%
\XVnote{\dag\ Of course there is nothing really special about the letters $x$
and $y$. Use of different letters, say $u$ and $v$, would give us a
\emph{uv-plane}.}

Given the coordinates of a point, if one marks that point in a figure, it is
then said that the point has been \emph{plotted}.

\begin{remark}
The point $O$ has coordinates $(0, 0)$. (Why?)
\end{remark}

If $P$ does not lie on either axis, neither coordinate of $P$ can be 0. If the
first coordinate of $P$ is 0, then $P$ is on the axis of ordinates. The
converse is also true.

\FIGXVSIX

\exercisegroup{15--7}
\begin{multicols}{2}
\begin{exlist}
\item Draw axes, and plot the points $(2, 4)$; $(-2, 4)$;
  $(-3, -\tfrac{1}{2})$; $(0, -5)$; $(-4, 0)$; and $(7\tfrac{1}{2}, -3)$.
\item A point $P$ lies in the first quadrant. The perpendicular distance from
  the $x$-axis is $4\tfrac{1}{2}$ units; the distance from the $y$-axis is 3
  units. What are the coordinates of $P$? What are the coordinates if
  distances are the same but $P$ is in the second quadrant?
\item[\stex 3.] Suppose that, using the same lines as axes, the unit of
  distance is chosen twice as large. What would be the effect upon the
  coordinates of points?
\item[\stex 4.] Where are all the points in the plane with abscissa 0?
  abscissa negative? ordinate positive?
\item[\stex 5.] Where are all the points in the plane whose ordinates are
  equal to their abscissas?
\item[6.] Draw the line that passes through the points $(-1, 5)$ and
  $(2, -1)$. Can you easily draw a line parallel to it, which passes through
  the origin?
\item[\stex 7.] Carry out the construction of coordinates when the axes are
  not perpendicular.
\item[8.] Both coordinates of the point $S$ are negative numbers. The points
  $S$ and $T$ are on the same side of the $x$-axis, but on opposite sides of
  the $y$-axis. What can you say about the coordinates of $T$?
% Manual \item[...] labels do not advance the list counter; resync so the
% plain items below number 9-17 instead of restarting at 3.
\setcounter{exlisti}{8}
\item The positive and negative $x$- and $y$-axes are rays from $O$. Describe
  the four quadrants as interiors of angles formed by these rays.
\item The abscissa of a point is 0. Where must it lie?
\item The ordinate of a point is 0. Where can it be?
\item What can you say about a point if the product of its coordinates is
  negative? 0? positive?
\item Plot five distinct points with first coordinate equal to 3. Describe the
  set of all points in the plane such that the first coordinate of every point
  of the set is 3.
\item Consider the set of all points in the plane such that the sum of the
  coordinates of each point is 10. Plot six such points, two in each of
  quadrants I, II, and IV. Why are there no such points in quadrant III?
\item Describe the set of all points $(x, y)$ such that $x + 1 < 0$.
\item What is the perpendicular distance from $(3, -4)$ to the $x$-axis? to
  the $y$-axis? from $(-5, -7)$ to each axis?
\item How far is $(3, 4)$ from $(0, 0)$?
\end{exlist}
\end{multicols}

%========================================================= 15-4
\section*{15--4\quad STRAIGHT LINES}
\addcontentsline{toc}{section}{15--4\quad Straight lines}

Having coordinates, we naturally ask how lines may be described by
coordinates. In exercises we have already examined a few lines.

For example, the set of all points $(x, y)$ with abscissa 3 is the line shown
in Fig.~15--7, and every point on the line parallel to the $y$-axis, at a
distance 3 to the right of the $y$-axis, has its $x$-coordinate equal to 3.
Hence, this line can be specified briefly by writing just $x = 3$. We say then
that ``$x = 3$'' is the \emph{equation} of this line.

\FIGXVSEVEN

Conversely, when we see the equation $x = 3$, we ask ourselves, ``What is the
set of all points $(x, y)$ for which $x = 3$?''

In the same way, all lines parallel to the $y$-axis have an equation $x = a$,
where $a$ is some real number, positive or negative. (If $a = 0$, what is the
line?)

Likewise, lines parallel to the $x$-axis will have equations $y = b$.

\exercisegroup{15--8}
\begin{multicols}{2}
\begin{exlist}
\item A line is parallel to the $y$-axis and passes through the point
  $(-3, 5)$. What is its equation? What if it is parallel to the $x$-axis?
\item What is the equation of the line that bisects the first and third
  quadrants? the second and fourth?
\item[\stex 3.] Plot several points $(x, y)$ whose coordinates satisfy the
  equation $x + y = 5$. Can you describe briefly the set of all points
  $(x, y)$ whose coordinates satisfy this equation?
\end{exlist}
\end{multicols}

We will prove shortly that every straight line has an equation of a simple
kind. But first we see that such simple equations turn out to be straight
lines.

\begin{XVexample}[ 1]
Consider the set of all points $(x, y)$ whose coordinates satisfy the equation
$2x - 3y + 6 = 0$. We first find a few of these points, by choosing a value
for $x$ (or $y$) and then solving for $y$ (or $x$). Thus, if $x = 0$, we have
$0 - 3y + 6 = 0$, or $y = 2$. If $x = 3$, then $y = 4$. If $y = 0$, then
$x = -3$.
\par\medskip
\begin{center}
\begin{tabular}{c|c}
\multicolumn{1}{c}{$x$} & \multicolumn{1}{c}{$y$} \\
0 & 2 \\
3 & 4 \\
$-3$ & 0 \\
\end{tabular}
\end{center}
\end{XVexample}

\FIGXVEIGHTNINE

When these points are plotted, we see (by looking at Fig.~15--8) that they
apparently lie on a line.

Find other points whose coordinates satisfy the equation, and plot them.

\begin{XVexample}[ 2]
Consider the set of all points $(x, y)$ whose coordinates satisfy the equation
$-4y - 8 = 0$. From the equation, we see that $y = -2$, regardless of what $x$
is. The set of all points $(x, -2)$ is the line parallel to the $x$-axis and a
distance 2 units below. See Fig.~15--9.
\end{XVexample}

\exercisegroup{15--9}
Plot several points $(x, y)$ whose coordinates satisfy the following
equations:
\begin{multicols}{2}
\begin{exlist}
\item $2x + 3y - 6 = 0$.
\item $x - 2y + 4 = 0$.
\item $3x - 9 = 0$.
\item $-x - y + 4 = 0$.
\item $2ax + 3ay - 6a = 0$; $a \neq 0$.
\item $-3x + 5y - 15 = 0$.
\item $y = 2x + 2$.
\item $x = 3y - 6$.
\end{exlist}
\end{multicols}

From the examples we observe that equations of the form $Ax + By + C = 0$,
where $A$, $B$, and $C$ are given numbers, are satisfied by the coordinates of
points on lines. We now prove this theorem.

\begin{XVstheorem}
If $A$, $B$, and $C$ are given numbers, with not both $A$ and $B$ zero, then
all points $(x, y)$ whose coordinates satisfy the equation $Ax + By + C = 0$
lie on a single line.
\end{XVstheorem}

\noindent\emph{Proof.}
We separate the proof into three cases.

\emph{Case} 1. $A = 0$. Then $B \neq 0$ (why?) and the equation is
$By + C = 0$, or $y = -C/B$. Then all points with this $y$-coordinate lie on
the line parallel to the $x$-axis and passing through $(0, -C/B)$.

\emph{Case} 2. $B = 0$. In this case, the points lie on a line parallel to the
$y$-axis.

\emph{Case} 3. $A \neq 0$ and $B \neq 0$. It suffices to show that any two
points $(x_1, y_1)$ and $(x_2, y_2)$ whose coordinates satisfy the equation
lie on one line with the point $(0, -C/B)$. This will be the case if and only
if the right triangles indicated in Fig.~15--10 are similar, where

\begin{equation}\tag{15--1}
y_1 = \frac{-Ax_1-C}{B}; \qquad y_2 = \frac{-Ax_2-C}{B}.
\end{equation}

\FIGXVTENELEVEN

But these triangles are similar if and only if

\begin{equation}\tag{15--2}
\frac{y_2-(-C/B)}{x_2} = \frac{y_1-(-C/B)}{x_1}.
\end{equation}

From Eq.~(15--1) we see that
\[
\frac{y_2+C/B}{x_2} = \frac{-A}{B} = \frac{y_1+C/B}{x_1}.
\]

Hence we conclude that the triangles are similar and the points
$[0, -(C/B)]$, $(x_1, y_1)$, and $(x_2, y_2)$ are collinear.

\begin{XVstheorem}
If $l$ is a line, then there are numbers $A$, $B$, and $C$, with not both $A$
and $B$ zero, such that every point $(x, y)$ on $l$ has coordinates that
satisfy the equation
\begin{equation}\tag{15--3}
Ax + By + C = 0.
\end{equation}
\end{XVstheorem}

\noindent\emph{Proof.}
As in the previous theorem, we consider three cases.

\emph{Case} 1. $l$ is parallel to the $x$-axis. Then there is a point
$(0, b)$ on $l$, and every point $(x, y)$ on $l$ satisfies $y = b$. This is of
the form of Eq.~(15--3), with $A = 0$, $B = 1$, and $C = -b$.

\emph{Case} 2. $l$ is parallel to the $y$-axis. This case is left to the
student.

\emph{Case} 3. $l$ is parallel to neither axis. Choose $(x_1, y_1)$ and
$(x_2, y_2)$, distinct points on $l$. Then $x_1 \neq x_2$ and
$y_1 \neq y_2$. (Why?) Let $(x, y)$ be any other point on $l$.

As shown in Fig.~15--11, the points $(x_2, y_1)$ and $(x, y_1)$ form, with
$(x_1, y_1)$, $(x_2, y_2)$, and $(x, y)$, two right triangles that are
similar. Because of similarity,\dag

\begin{equation}\tag{15--4}
\frac{y-y_1}{y_2-y_1} = \frac{x-x_1}{x_2-x_1}
\end{equation}
\XVnote{\dag\ Figure 15--11 is drawn so that $x_1 < x_2 < x$, and
$y_1 < y_2 < y$, but all orders can occur for different points and different
lines. In every case, however, Eq.~(15--4) will be satisfied. Thus it makes no
difference which point is considered as $(x_1, y_1)$, so long as one is
consistent in its use.}

\noindent or
\[
(y_2 - y_1)x - (x_2 - x_1)y + x_2y_1 - y_2x_1 = 0,
\]
which is of the form of Eq.~(15--3). This completes the proof.

\begin{XVremark}{1}
An equation equivalent to Eq.~(15--4) is
\begin{equation}\tag{15--5}
y - y_1 = \frac{y_2-y_1}{x_2-x_1}\,(x-x_1).
\end{equation}
This is a useful equation to remember because it enables one to write at once
an equation satisfied by the coordinates of any point on a line through two
given points.
\end{XVremark}

\begin{XVremark}{2}
Two different equations may correspond to the same line. Thus for example,
$x = 2$ and $2x - 4 = 0$ represent the same line. Such equations are said to
be \emph{equivalent} equations. In arriving at equations for lines, we usually
do not care which equivalent equation we obtain.
\end{XVremark}

\begin{definition}
If $Ax + By + C = 0$ is an equation satisfied by the coordinates $(x, y)$ of
all points on a line $l$, we say that $Ax + By + C = 0$ is the equation of
$l$. We also speak of the \emph{``line $Ax + By + C = 0$.''} Such an equation
of the first degree in $x$ and $y$ will be called a \emph{linear equation} in
$x$ and $y$.
\end{definition}

\FIGXVTWELVE

\begin{XVexample}
Find the equation of the line through $(-2, 3)$ and $(1, 5)$ in Fig.~15--12,
and draw the line. Using Eq.~(15--5), we have
\[
y - 3 = \frac{5-3}{1-(-2)}\,[x-(-2)]
\qquad\text{or}\qquad
y = \tfrac{2}{3}x + \tfrac{13}{3}.
\]
\end{XVexample}

\exercisegroup{15--10}
\begin{multicols}{2}
\begin{exlist}
\item Draw the lines
  \begin{enumerate}[label=(\alph*),leftmargin=1.5em,itemsep=0pt,topsep=2pt,
    parsep=0pt]
  \item $3x - 4y = 12$,
  \item $3x - 4y = 8$,
  \item $4x + 3y = 0$,
  \item $4x + 15 = 0$,
  \item $-2x - 3y + 5 = 0$,
  \item $y = \dfrac{-1}{2}\,x + 3$,
  \item $y = \dfrac{-1}{2}\,x$,
  \item $x = 3y + 3$,
  \item[(k)] $5y - 12 = 0$,
  \item[(l)] $2x + y = 7$.
  \end{enumerate}
\item Find the equations of the lines through the following pairs of points:
  \begin{enumerate}[label=(\alph*),leftmargin=1.5em,itemsep=0pt,topsep=2pt,
    parsep=0pt]
  \item $(-1, 1)$ and $(3, -1)$,
  \item $(2, 3)$ and $(-2, -3)$,
  \item $(a, 0)$ and $(0, b)$,
  \item $(-3, 0)$ and $(0, -5)$,
  \item $(-2, -4)$ and $(-5, -5)$,
  \item $(0, b)$ and $(1, m + b)$.
  \end{enumerate}
\item In Exercise 1, how many points did you plot to draw each line? Is there
  a reason for plotting more than two?
\item Find the coordinates of the point of intersection of the lines
  $3x - 5y = 7$ and $2x + 3y = -1$.
\item A line meets the $y$-axis in $(0, -3)$ and the $x$-axis in $(-4, 0)$.
  Where does this line intersect the line $y = 7x - 19$?
\item Show that the lines $2x + 3y = 0$ and $4x + 6y = 5$ are parallel.
\end{exlist}
\end{multicols}

%========================================================= 15-5
\section*{15--5\quad DISTANCE}
\addcontentsline{toc}{section}{15--5\quad Distance}

The distance between any two points of the plane is easy to calculate from the
coordinates of the points. It is particularly easy if the line through the
points is parallel to one of the axes. Thus, $\sg{AB} = 1 - (-3) = 4$, and
$\sg{CD} = |x_2 - x_1| = |x_1 - x_2|$ in Fig.~15--13.

\FIGXVTHIRTEENFOURTEEN

If the line through the points is not parallel to one of the axes, then with
the point $Q$ $(x_2, y_1)$ a right triangle $P_1P_2Q$ is formed, as in
Fig.~15--14, and the Pythagorean theorem asserts that
$\sg{P_1P_2}^2 = \sg{P_1Q}^2 + \sg{P_2Q}^2 = |x_2 - x_1|^2 +
|y_2 - y_1|^2$. Taking square roots, we obtain
$\sg{P_1P_2} = \sqrt{(x_2-x_1)^2 + (y_2-y_1)^2}$.

\begin{remark}
In using this formula it makes no difference which point is considered to be
$(x_1, y_1)$ because $(x_1 - x_2)^2 = (x_2 - x_1)^2$ and
$(y_1 - y_2)^2 = (y_2 - y_1)^2$.
\end{remark}

\exercisegroup{15--11}
\begin{multicols}{2}
\begin{exlist}
\item Show that the triangle with vertices $(2, 6)$, $(5, 7)$, and
  $(-8, -2)$ is a right triangle.
\item What kind of triangle is the one with vertices $(3, 2)$, $(-4, 4)$, and
  $(1, -5)$?
\item[\stex 3.] Write an equation that states that the point $(x, y)$ is
  equidistant from $(-3, 2)$ and $(2, 1)$. Show that this equation is
  equivalent to a linear equation. Hence conclude that the set of all points
  equidistant from $(-3, 2)$ and $(2, 1)$ is a straight line.
\item[4.] Show that $(-1, -18)$, $(-13, -2)$, $(11, -8)$, and $(-1, 8)$ are
  vertices of a parallelogram.
% Manual \item[...] labels do not advance the list counter; resync so the
% plain items below number 5-7 instead of restarting at 3.
\setcounter{exlisti}{4}
\item Show that if $(x, y)$ is a point on a circle whose center is $(-2, 1)$
  and radius is 3, then $x^2 + y^2 + 4x - 2y = 4$.
\item Show that the points $(2, 1)$, $(3, \tfrac{1}{2})$, and $(-2, 3)$ lie on
  a line. What geometric theorems do you use?
\item Show that the distance formula yields $|x_1 - x_2|$ if $y_1 = y_2$.
\end{exlist}
\end{multicols}

%========================================================= 15-6
\section*{15--6\quad MIDPOINT}
\addcontentsline{toc}{section}{15--6\quad Midpoint}

The midpoint of the segment $P_1P_2$ may be shown to have coordinates
\[
\left(\frac{x_1+x_2}{2},\ \frac{y_1+y_2}{2}\right).
\]

The proof is left to the student. If the line through the points is parallel
to one of the axes, the proof is very simple. If $x_1 \neq x_2$ and
$y_1 \neq y_2$, a right triangle may be formed as in Fig.~15--15, and
parallels to the coordinate axes through the midpoint will bisect the sides in
points whose coordinates are easily computed.

\begin{XVexample}
Show in Fig.~15--16 that the midpoint formula actually gives a point
equidistant from $(-1, 3)$ and $(1, -1)$.
\end{XVexample}

\FIGXVFIFTEENSIXTEEN

The point $M$ is
$\left(\dfrac{-1+1}{2},\ \dfrac{3+(-1)}{2}\right) = (0, 1)$. Then
\[
\sg{P_1P_2} = \sqrt{2^2 + 4^2} = \sqrt{20} = 2\sqrt{5},
\]
\[
\sg{P_1M} = \sqrt{1^2 + 2^2} = \sqrt{5},
\]
and
\[
\sg{MP_2} = \sqrt{1^2 + 2^2} = \sqrt{5}.
\]
Hence
\[
\sg{P_1P_2} = 2\sqrt{5} = \sqrt{5} + \sqrt{5} = \sg{P_1M} + \sg{MP_2}.
\]

\exercisegroup{15--12}
\begin{multicols}{2}
\begin{exlist}
\item Find the midpoints of the sides, and the lengths of the medians, of the
  triangle whose vertices are $(-4, 10)$, $(12, 10)$, and $(4, -2)$.
\item Find the coordinates of the points that divide the segment $(-3, -7)$
  $(11, 17)$ into fourths.
\item[\stex 3.] Find the points that trisect the segment $(-3, -7)$
  $(6, 11)$.
\item[\stex 4.] Derive a formula for the coordinates of the points that
  trisect the segment $(x_1, y_1)$ $(x_2, y_2)$.
\item[\stex 5.] Show that the medians of the triangle with vertices
  $(x_1, y_1)$, $(x_2, y_2)$, and $(x_3, y_3)$ intersect in the point
  $[(x_1 + x_2 + x_3)/3, (y_1 + y_2 + y_3)/3]$.
\item[6.] The point $(-3, 4)$ is the midpoint of a segment, one end of which
  is $(5, 7)$. Find the other end.
\item[\stex 7.] In Fig.~15--17, show that the line joining the midpoints of
  the sides of a triangle is half as long as the third side, and parallel to
  it. [\emph{Hint:} Choose your coordinate system so that one vertex is at the
  origin and another is on the positive $x$-axis.]
\item[\stex 8.] Prove that the diagonals of a parallelogram bisect each
  other.
\end{exlist}
\end{multicols}

\FIGXVSEVENTEEN

%========================================================= 15-7
\section*{15--7\quad CURVES AND EQUATIONS}
\addcontentsline{toc}{section}{15--7\quad Curves and equations}

We have seen that certain sets of points can be described by equations. Thus a
straight line has a linear or first degree equation, in the sense that the
coordinates of every point on the line satisfy an equation of first degree in
$x$ and $y$ of the form $Ax + By + C = 0$, and conversely. In the same way,
other sets of points can be described by different equations. Perhaps the
simplest example of this is furnished by circles.

\FIGXVEIGHTEENNINETEEN

Suppose that $(x, y)$ in Fig.~15--18 is a point on the circle whose center is
$(h, k)$ and whose radius is $r$, then
\[
\sqrt{(x-h)^2 + (y-k)^2} = r
\]
or, equivalently,

\begin{equation}\tag{15--6}
(x-h)^2 + (y-k)^2 = r^2.
\end{equation}

Conversely, a point $(x, y)$ whose coordinates satisfy Eq.~(15--6) will be at
a distance $r$ from $(h, k)$ and therefore lies on the circle.

We speak of Eq.~(15--6) as ``the'' equation of a circle, although any
equivalent equation would also be called ``the'' equation. For example, the
equation of the circle whose center is at $(-2, 3)$, and whose radius is 4, is
\[
(x + 2)^2 + (y - 3)^2 = 16
\]
or
\[
x^2 + y^2 + 4x - 6y = 3.
\]

Given an equation in the form of Eq.~(15--6), we can recognize it at once for
the equation of a circle. Any equation equivalent to Eq.~(15--6) is the
equation of a circle. For example, consider the equation
\[
x^2 + y^2 + 6x - 8y + 19 = 0.
\]

Rewrite the equation as
\[
(x^2 + 6x + \,?) + (y^2 - 8y + \,?) = -19 + \,? + \,?.
\]

By completing the squares of the expressions in parentheses, we have
\[
(x^2 + 6x + 9) + (y^2 - 8y + 16) = -19 + 9 + 16
\]
or
\[
(x + 3)^2 + (y - 4)^2 = 6,
\]
which is the equation of a circle whose center is at $(-3, 4)$ and whose
radius is $\sqrt{6}$. See Fig.~15--19.

Other curves (sets of points) will have other kinds of equations. The student
is referred to books on analytic geometry for more detail.

\exercisegroup{15--13}
\begin{multicols}{2}
\begin{exlist}
\item Write the equation of the circle whose center is $(3, 5)$ and whose
  radius is 6.
\item Write the equation of the circle whose center is at $(-4, -3)$ and which
  (a) is tangent to the $y$-axis, (b) passes through the origin.
\item Find center and radius of the following circles:
  \begin{enumerate}[label=(\alph*),leftmargin=1.5em,itemsep=0pt,topsep=2pt,
    parsep=0pt]
  \item $x^2 + y^2 + 8x - 9 = 0$,
  \item $x^2 + y^2 - 6x - 10y = 2$,
  \item $x^2 + y^2 - 7 = 0$,
  \item $2x^2 + 2y^2 + 5x + 4y - 5 = 0$,
  \item $x^2 + y^2 - 6x - 8y + 25 = 0$.
  \end{enumerate}
\item[\stex 4.] Find an equation satisfied by the coordinates of every point
  $(x, y)$ whose distance from $(2, 0)$ is equal to its distance from the line
  whose equation is $x + 2 = 0$.
\item[\stex 5.] If the sum of the distances from $(x, y)$ to $(-3, 0)$ and
  from $(x, y)$ to $(3, 0)$ is equal to 10, show that
  $16x^2 + 25y^2 = 400$.
\end{exlist}
\end{multicols}

\end{document}
